\section{Partitions and the law of total probability}

Conditional probability becomes especially useful when a problem can be split
into cases. Many probability questions have the following structure: we want the
probability of an event \(A\), but the mechanism that produces \(A\) depends on
which one of several possible background cases occurred.

The mathematical language for such a situation is a partition.

\subsection*{Splitting the sample space into cases}

Suppose the sample space \(\Omega\) is split into events
\[
B_1,B_2,\ldots,B_n.
\]
These events should represent mutually exclusive and exhaustive cases. Mutually
exclusive means that no two cases can occur together:
\[
B_i\cap B_j=\varnothing\qquad\text{when }i\neq j.
\]
Exhaustive means that every possible outcome belongs to one of the cases:
\[
B_1\cup B_2\cup\cdots\cup B_n=\Omega.
\]

\begin{definitionbox}
A finite collection of events \(B_1,
\ldots,B_n\) is a \textbf{partition} of
\(\Omega\) if the events are pairwise disjoint and their union is \(\Omega\). In
applications of conditional probability, we usually also require
\[
P(B_i)>0
\]
for each case that appears in a conditional probability.
\end{definitionbox}

A partition is a disciplined way to say: exactly one of these cases occurs.

\begin{examplebox}
If a person is selected from a population, the events
\[
B_1=\{\text{person is in group 1}\},\quad
B_2=\{\text{person is in group 2}\},\quad
B_3=\{\text{person is in group 3}\}
\]
form a partition if every person belongs to exactly one of the three groups.
\end{examplebox}

\subsection*{How an event is split by a partition}

Let \(A\) be any event. If \(B_1,
\ldots,B_n\) is a partition of \(\Omega\), then
\(A\) can be split according to the cases:
\[
A=(A\cap B_1)\cup(A\cap B_2)\cup\cdots\cup(A\cap B_n).
\]
These pieces are disjoint because the \(B_i\)'s are disjoint. Thus the
probability of \(A\) is the sum of the probabilities of the pieces:
\[
P(A)=P(A\cap B_1)+P(A\cap B_2)+\cdots+P(A\cap B_n).
\]

This identity is just additivity applied after decomposing \(A\) into disjoint
parts.

\subsection*{Replacing intersections by conditional probabilities}

For each case \(B_i\), the multiplication rule gives
\[
P(A\cap B_i)=P(A\mid B_i)P(B_i).
\]
Substituting this into the previous sum gives
\[
P(A)=P(A\mid B_1)P(B_1)+P(A\mid B_2)P(B_2)+\cdots+P(A\mid B_n)P(B_n).
\]
Equivalently,
\[
P(A)=\sum_{i=1}^n P(A\mid B_i)P(B_i).
\]

\begin{theorembox}
\textbf{Law of total probability.} If \(B_1,
\ldots,B_n\) is a partition of
\(\Omega\) and \(P(B_i)>0\), then for every event \(A\),
\[
P(A)=\sum_{i=1}^n P(A\mid B_i)P(B_i).
\]
\end{theorembox}

The formula says that the total probability of \(A\) is obtained by looking at
all possible cases. In each case, we multiply the probability of the case by the
probability of \(A\) inside that case, and then add the contributions.

\begin{remarkbox}
The law of total probability is not a new mysterious rule. It is additivity over
a partition, with each piece \(P(A\cap B_i)\) rewritten as
\(P(A\mid B_i)P(B_i)\).
\end{remarkbox}

\subsection*{Weighted averages of conditional probabilities}

The expression
\[
P(A)=\sum_{i=1}^n P(A\mid B_i)P(B_i)
\]
can be read as a weighted average of the conditional probabilities
\[
P(A\mid B_1),\ldots,P(A\mid B_n).
\]
The weights are
\[
P(B_1),\ldots,P(B_n).
\]
They are non-negative and add to one because the \(B_i\)'s form a partition:
\[
P(B_1)+\cdots+P(B_n)=1.
\]

This interpretation is useful. If \(A\) is likely in a case that rarely occurs,
its contribution to the total probability may still be small. If \(A\) is only
moderately likely in a case that occurs often, its contribution may be large.

\subsection*{Example: choosing one of two coins}

Suppose a box contains two coins. Coin 1 is fair:
\[
P(H\mid B_1)=\frac12.
\]
Coin 2 is biased toward heads:
\[
P(H\mid B_2)=\frac34.
\]
A coin is chosen at random, with
\[
P(B_1)=P(B_2)=\frac12,
\]
and then the chosen coin is tossed once. What is the probability of heads?

Let
\[
A=\{\text{heads occurs}\}.
\]
The events \(B_1\) and \(B_2\) form a partition of the background mechanism:
we either chose coin 1 or coin 2. The law of total probability gives
\[
P(A)=P(A\mid B_1)P(B_1)+P(A\mid B_2)P(B_2).
\]
Thus
\[
P(A)=\frac12\cdot\frac12+\frac34\cdot\frac12
=\frac14+\frac38=\frac58.
\]

The answer is between \(1/2\) and \(3/4\), because before the toss we do not know
which coin was chosen. The total probability is the average of the head
probabilities of the two coins, weighted by the probabilities of choosing them.

\subsection*{Example: urn chosen before drawing}

Suppose one urn is selected first, and then one object is drawn from that urn.
Urn 1 contains \(3\) red objects and \(1\) blue object. Urn 2 contains \(1\) red
object and \(4\) blue objects. The probability of choosing urn 1 is \(2/3\), and
the probability of choosing urn 2 is \(1/3\).

Let
\[
A=\{\text{the drawn object is red}\},
\]
\[
B_1=\{\text{urn 1 was chosen}\},\qquad
B_2=\{\text{urn 2 was chosen}\}.
\]
Then
\[
P(A\mid B_1)=\frac34,\qquad P(A\mid B_2)=\frac15.
\]
The cases \(B_1\) and \(B_2\) form a partition. Therefore
\[
P(A)=P(A\mid B_1)P(B_1)+P(A\mid B_2)P(B_2).
\]
Substituting the values,
\[
P(A)=\frac34\cdot\frac23+\frac15\cdot\frac13.
\]
Hence
\[
P(A)=\frac12+\frac1{15}=\frac{17}{30}.
\]

The result is not obtained by averaging the contents of the two urns equally,
because the urns are not chosen equally. The weights \(2/3\) and \(1/3\) matter.

\subsection*{Example: data table}

Suppose a population consists of two departments. Department 1 contains \(40\%\)
of the population, and department 2 contains \(60\%\). In department 1, \(70\%\)
of people use a certain system. In department 2, \(30\%\) use it.

Let
\[
A=\{\text{the selected person uses the system}\},
\]
\[
B_1=\{\text{the selected person is in department 1}\},
\]
\[
B_2=\{\text{the selected person is in department 2}\}.
\]
Then
\[
P(B_1)=0.4,\qquad P(B_2)=0.6,
\]
\[
P(A\mid B_1)=0.7,\qquad P(A\mid B_2)=0.3.
\]
Therefore
\[
P(A)=0.7\cdot0.4+0.3\cdot0.6=0.28+0.18=0.46.
\]

The overall usage rate is \(46\%\). It is a weighted average of the department
usage rates, weighted by department sizes.

\subsection*{Choosing the right partition}

The law of total probability is useful only after the cases have been chosen
well. A good partition should satisfy three conditions.

First, the cases must be mutually exclusive. If two cases can occur together,
then adding their contributions may double-count outcomes.

Second, the cases must be exhaustive. If some outcomes are not covered by any
case, then the sum will miss part of the probability.

Third, the conditional probabilities inside the cases should be meaningful and
usually easier to understand than the original probability.

\begin{examplebox}
For a test problem, a natural partition may be
\[
B_1=\{\text{person has the condition}\},\qquad
B_2=\{\text{person does not have the condition}\}.
\]
These two events are disjoint and exhaustive. They also match how test accuracy
is usually described: sensitivity gives \(P(\text{positive}\mid B_1)\), and
false positive rate gives \(P(\text{positive}\mid B_2)\).
\end{examplebox}

\subsection*{Common mistake: forgetting the weights}

A common mistake is to average conditional probabilities without using the case
probabilities. For example, if
\[
P(A\mid B_1)=0.9,\qquad P(A\mid B_2)=0.1,
\]
one might be tempted to say that \(P(A)=0.5\). This is correct only if
\[
P(B_1)=P(B_2)=0.5.
\]
If instead
\[
P(B_1)=0.1,\qquad P(B_2)=0.9,
\]
then
\[
P(A)=0.9\cdot0.1+0.1\cdot0.9=0.18.
\]
The rare high-probability case contributes less than the common low-probability
case.

\subsection*{What the law of total probability teaches}

The law of total probability teaches how to assemble a probability from cases.
The event \(A\) is decomposed into disjoint pieces \(A\cap B_i\). Each piece is
measured by first entering case \(B_i\), then measuring \(A\) within that case.

\begin{theorembox}
To use total probability, ask:
\[
\text{Which mutually exclusive and exhaustive cases can produce the event?}
\]
Then add the case contributions:
\[
\text{probability of case}\times\text{probability of event inside case}.
\]
\end{theorembox}

This prepares the next idea. If total probability lets us compute the probability
of an observation from possible causes or cases, Bayes' theorem lets us reverse
the question: after the observation is known, how likely is each case?
